\section{Numerical cross-checks}
\label{app:numerics}

The closed forms were compared with computations that do not use the steps of the proofs. Solvers: HiGHS \cite{highs} through SciPy \cite{scipy} for linear programmes, CLARABEL \cite{clarabel} through CVXPY \cite{cvxpy} for the semidefinite programme.

\emph{Support function of $\cR$ (route V6).} For $24$ directions $(a,b)$ and $16$ pure states $\psi_\alpha$ the maximum of $am_1+bm_2$ over $\cQ$ was computed as a semidefinite programme in the primal form of Lemma~\ref{lem:sos}(b). The maximum over $\alpha$ differs from Eq.~\eqref{eq:hR} by at most $\sdpDev$.

\emph{Linear programmes without the covariant reduction (routes V6, P4, P5).} For a direction in the plane of the ratios, the largest $t$ was computed for which $\Lambda_{t\hat\eta}=\sum_k|Z_k\rangle\!\rangle\langle\!\langle G_k|$ with a fixed finite family of operators $G_k$ (pure states for $\cR$; operators of the form \eqref{eq:squares} for $\cR_2$) and Hermitian variables $Z_k$ whose Husimi functions are nonnegative on a grid of directions. The finite family makes the programme an inner approximation and the finite grid an outer relaxation, so the two discretisations act with opposite signs. For $\cR$, with $\lpRK$ operators and $\lpRnq$ grid directions, the ratio of the computed radius to the closed form lies in $[\lpRlo,\lpRhi]$ over $\lpRdirs$ directions; for $\cR_2$, with $\lpRtwoK$ operators and $\lpRtwonq$ grid directions, in $[\lpRtwolo,\lpRtwohi]$ over $\lpRtwodirs$ directions. The three largest deviations were traced to the discretisation by refinement: at the top of $\cR$ the ratio moves from $\refA$ to $\refB$ when the grid is refined from $\refAnq$ to $\refBnq$ directions; at the apex of $\cR_2$ from $\refC$ to $\refD$ (grid $\refCnq$ to $\refDnq$); and at $80^\circ$ in $\cR_2$, where the boundary point $s=1.81$ fell between members of the operator family, from $\refE$ to $\refF$ after adding operators with $\tau=1.8$. These runs support the closed forms at the level of a few parts in a thousand; they are not proofs.

\emph{Two spin-$1/2$ particles (route H10).} The threshold of Theorem~\ref{thm:pair} was bracketed by two linear programmes: on finite subsets of the steered states and of the effects, where infeasibility implies contextuality of the full fragment, and on finite subsets of the two dual cones, where feasibility implies noncontextuality. For a random unpolarised family $c=tc_0$ with $\|c_0\|_1=1$ and for a polarised family with $\kappa_A=0.75$, $\kappa_B=1$, the exact threshold lies inside the bracket in all runs, and the bracket width is $\bracketWidth$ for $\bracketN$ points. For the $\Lambda\bar\Lambda$ coefficients the first programme is infeasible at $\cos\theta=0$ and $0.4$ and the second is feasible at $\cos\theta=0.65$, in agreement with the window $|\cos\theta|<\cosc$.

\emph{Spin $j$ (routes DV8, G4).} For $j\le4$ the linear programme that maximises $\Tr[ZJ_z]/(j+1)$ over diagonal unit-trace $Z$ with nonnegative Husimi function on a grid of $2\times10^4$ polar angles returns $\xi_{2j}$ to $\gjDev$ (Table~\ref{tab:gj}), and the kernels built from the optimisers are probability vectors with $\kappa_1=\xi_{2j}$ and nonzero multipole ratios.

\emph{First-moment bound (route G7).} For degrees $n=1,\dots,8$ the normalised first moment of random nonnegative, non-zonal polynomials on the sphere (sums of squares, multiplied by $1+\bm{m}\!\cdot\!\bq$ for odd $n$) stays below $\xi_n$, and the extremal zonal polynomial of Appendix~\ref{app:spinj} attains $\xi_n$ to rounding accuracy.

\emph{Lemmas (routes V2, H12).} The constructions of Lemma~\ref{lem:sos}(a) and Lemma~\ref{lem:range} were run on $200$ random instances each, with residuals below $10^{-9}$.
